\documentclass[11pt]{article}
\usepackage[margin=1in]{geometry}
\usepackage{amsmath,amssymb,amsthm}
\usepackage{booktabs}
\usepackage{graphicx}
\usepackage{natbib}
\usepackage{hyperref}
\newtheorem{theorem}{Theorem}
\newtheorem{lemma}{Lemma}
\newtheorem{proposition}{Proposition}
\newtheorem{corollary}{Corollary}
\newtheorem{conjecture}{Conjecture}
\newtheorem{definition}{Definition}
\newcommand{\sd}{\mathrm{sd}}
\title{A memoryless two-header tape machine:\\
what linear time explains, and what our predictor costs}
\author{}
\date{}

\begin{document}
\maketitle

\begin{abstract}
A tape of $n$ cells, each the average of a pair of bits, is read by two headers that move inside a
region. At each step the rule evaluates a fixed function of the window between them and chooses one of
eight actions --- move a header, cut a cell from the left, the right or both, or halt --- and outputs
the position where the headers meet. We ask what linear time explains, and what a predictor with little
memory recovers.

For the canonical rule we construct an exact simulator costing $5.3$ reads per cell, flat from $n=32$
to $1024$; the same construction covers all twenty-four move tables and steps of one to three cells
at $n=256$, and the two rules that first looked expensive are not in that class either. For four further rules no such simulator is known: their
trajectories are quadratic (a fitted exponent $2.089\pm.026$ on the median and $1.977\pm.042$ on the
mean), and a simulation carrying $O(\log n)$ bits costs $0.62\,n^2$ reads, about eleven reads per step. On the first of them, a predictor confined to $o(n^2)$ reads recovers $R^2=.58$ of the output at $n=256$, with two seeds agreeing to $.004$, and so any ceiling constant for these rules is at least $.56$ once the selection
over budgets is accounted for --- at $n=256$; at $n=512$ the same procedure gives $.37$, and we
claim no mode of convergence, so this is a finite-$n$ lower bound and not an asymptotic one.

For the expensive side we prove a chain of structural facts, each verified by exhaustive computation: the second bit is a
prefix-average comparison; with the other bit forced false the rule is exactly a largest-concavity
scan, checked on every admissible sequence for $n=6,\dots,13$. A counting
argument over the cells a machine reads cannot separate the classes --- such a bound can never exceed $n$, and $n$ is
on the cheap side --- and neither can an argument about memory, since the rule is simulable with
$\Theta(\log n)$ bits. What separates them is repeated reads: the scan costs $\Theta(n^2)$ reads in
expectation, and a single-head machine with $O(\log n)$ bits needs $\Omega(n^2/\log n)$ of them once
the cut-to-cut communication is at its measured size, so no such machine decides the output in linear
time.

\end{abstract}

\section{Introduction}

Two papers precede this one. The first proposes an arcsine limit for a triangular walk on a random
binary matrix \citep{luo-tww}; the second varies the specification and maps the shapes that result
\citep{luo-family}. The present paper studies the machine behind that walk without its
cycle-handling mechanism, and it is organised around two questions.

\emph{How much does linear time explain?} For the rules of this class that we can put in it, all of
it: if the moves are bounded and the cuts fall on the boundary, the output is computable exactly in
$O(n)$ reads with $O(\log n)$ bits of state, so no fitted model is needed and no ceiling applies. Four
others are candidates for the expensive class: their bits are cheap and their trajectories are long,
and they are the subject of Section~\ref{sec:walk} and of the ceiling in
Section~\ref{sec:ceiling}. Between the two sit four rules --- the interior cut, the far end, the
median and the rule that ignores the window --- whose simulators are linear per step but whose read
ratio grows with $n$ on the samples we have; we report their numbers and put none of them in either
class.

\emph{What does exact prediction cost?} Less than a rebuild suggests. A cut removes a boundary
cell, so the accumulators kept over the region lose one term; the window's are copied from them
when the headers restart, and nothing is rebuilt from the tape. A predictor that carries them follows the walk in
$O(n)$ reads with $O(\log n)$ bits of state, measured at $5.3\,n$ reads per tape from $n=32$ to
$1024$ (Table~\ref{tab:sim}), against $n$ to $2n$ for the machine itself, which is handed the two-bit
window function and never computes it.

\emph{What we can prove.} One direction only: if a rule's every move shrinks the window, and its
feature can be updated when one cell leaves, then the output is computable in $2n$ reads. The
converse is false, and Section~\ref{sec:barrier} explains why every natural attempt to state it dies
on a legal counterexample. In the other direction Section~\ref{sec:lower} proves two structural facts about the longest-run
rule and derives a corollary from them: it is a largest-concavity scan once one bit is forced
false, checked on every admissible sequence for $n=6,\dots,13$; and the quadratic cost of that scan
is an upper bound rather than a measurement. The consequence is negative and it is the point:
no argument that counts reads can separate the classes, because a linear budget already sits on the
cheap side. What has to be proved is a bound on memory.

\section{The machine}\label{sec:machine}

\paragraph{Tape and measure.} $\tau\in\{0,\tfrac12,1\}^n$: $n$ cells, the $i$-th the average of bit $i$ of column one and bit $i$ of column two.
Each of the two columns is drawn uniformly from the $\binom nh$ subsets of size $h=\lceil n/2\rceil$,
independently across columns, so for even $n$ the cell law is $P(0)=P(1)=\tfrac14$ and
$P(\tfrac12)=\tfrac12$, and for odd $n$, where $h=(n+1)/2$, it is
$\big(\frac{n-1}{2n}\big)^2$, $\frac{n^2-1}{2n^2}$ and $\big(\frac{n+1}{2n}\big)^2$; the cells are
negatively correlated. Tapes are drawn independently of one another.

\paragraph{Reads.} A \emph{read} is one access to one cell of the tape. The predictors below are
charged one read per cell they touch, so a pass over a window of length $m$ costs $m$ reads.

\paragraph{State and actions.} Fix $n\ge3$. A region $[L,R]$ and two headers $(\ell,r)$ with
$L\le\ell\le r\le R$, initially $L=\ell=0$ and $R=r=n-1$. At each step the rule chooses one of eight
actions:
\[
\ell\mapsto\ell+1,\qquad \ell\mapsto\ell-1,\qquad r\mapsto r+1,\qquad r\mapsto r-1,
\]
\[
L\mapsto L+1\ \ (\text{left cut}),\qquad R\mapsto R-1\ \ (\text{right cut}),\qquad
\text{both (double cut)},\qquad \text{halt}.
\]
The three cuts are actions in their own right. They may be taken at any step, not only
when a header has reached a boundary, and a rule may cut with both headers far from the ends.
Each cut removes one cell from each side it names, so the region only shrinks; after any cut the
headers restart at $(\ell,r)=(L,R)$. The tape is read-only.

\paragraph{Definition.} A \emph{strategy} is a fixed function
\[
f(\varphi(\tau[\ell..r]),\ n,\ \ell,\ r,\ L,\ R)\ \longrightarrow\
\{\ell{+}1,\ \ell{-}1,\ r{+}1,\ r{-}1,\ \text{left cut},\ \text{right cut},\ \text{double cut},\ \text{halt}\},
\]
where $\varphi$ is a fixed function of the window with a finite value set. The geometry is part of the
input: a strategy may depend on $\ell$, $r$, $L$ and $R$ as well as on the window. It has no other
memory, and in particular it cannot count its own steps.

A run is defined as follows.
\begin{enumerate}
\item Start at $(L,R,\ell,r)=(0,n-1,0,n-1)$.
\item If $\ell=r$, the run stops and outputs $\delta=\ell/(n-1)$.
\item If the state $(L,R,\ell,r)$ has occurred before, the run ends with no output: a rule without
memory repeats from there, and we stop rather than follow the repetition.
\item Otherwise the rule reads the window, evaluates $f$, and acts:
\begin{enumerate}
\item a move is legal only if it keeps both headers inside $[L,R]$, and if it brings them together or
past each other the run stops and outputs $\delta$ at the position the moving header reaches;
\item a cut is legal only if the region keeps at least one cell, and it acts on the region as
described above;
\item halt ends the run with no output.
\end{enumerate}
\end{enumerate}

\paragraph{Moving outside the region is not an action.} Forcing a header past a boundary is not a legal
action, and the model does not repair it: a rule that wants to discard the cell behind a boundary has
to say so, and the action that says so is a cut. The rules studied below are written as four-move
tables; read as strategies, they choose the cut on the side a move would have left and halt once the
region is small. Those are choices of those rules, not clauses of the model. A rule that halts
instead, or that keeps walking in a two-cell region, is an equally legal strategy and a different one.

\paragraph{Outcomes.} Three endings. A meeting outputs $\delta=\ell/(n-1)$. A halt outputs nothing. A
repeated state outputs nothing.

\section{Two classes of strategy}\label{sec:classes}

A strategy of this model is a window function $\varphi$ together with a decision function $f$ that
sees $\varphi$, both header positions and both boundaries, and returns one of the eight
actions of Section~\ref{sec:machine}. What separates strategies for our purposes is not
what they read but what their output costs.

\begin{definition}[Pointer and walk]
A strategy is a \emph{pointer} strategy if its output can be computed in $O(n)$ cell reads with
$O(\log n)$ bits of state --- that is, if we possess such a simulator. We use \emph{walk} for the
rest, and the word records the absence of a simulator we know rather than a proof that none exists:
Section~\ref{sec:lower} says what a proof would need.
\end{definition}

The classification is operational, and every claim below is a construction (a simulator with a read
count) unless it says otherwise. The second half of the definition is negative, so a rule is placed
on the expensive side only with the evidence we have: a simulator of cost $\Theta(n^2)$ and no
$O(n)$ one. We write "no $O(n)$ simulator is known" where that is all we can say, and
Section~\ref{sec:ceiling} states what would be needed to turn it into a lower bound.

\subsection{Simulators, and what they do and do not classify}
\label{sec:sims}

\paragraph{The lag-one autocorrelation.} One window function recurs below, in the canonical rule
and in the rule below that ignores the window: the lag-one autocorrelation of
a tape,
\[
\mathrm{ac}_1=\frac{\sum_{i=1}^{n-1}(a_i-\bar a)(a_{i+1}-\bar a)}{\sum_{i=1}^{n}(a_i-\bar a)^2},
\qquad \bar a=\frac1n\sum_{i=1}^{n}a_i,
\]
with $\mathrm{ac}_1=0$ when the denominator vanishes. It is computable in one pass with five running
scalars.

\paragraph{The canonical rule.} It is the special case in which $\varphi$ takes values in
$\{0,1\}^{2}$ and $f$ depends on the two bits alone, together with two conventions for what to do when the
table's move is illegal: it chooses the cut on the side the move would have left --- and the double
cut when both headers are on the boundary, which is the only case in which both sides are left at
once --- and it halts once the region has fewer than three cells. The double-cut case is what makes
Theorem~\ref{thm:concavity} read as it does. The two bits are $b_1$, the sign of the lag-one
autocorrelation of the window (with $b_1=[a_r>a_\ell]$ at length two and $[a_\ell>1/2]$ at length one),
and $b_2$, the comparison of the window mean with the mean of its first
$\max(1,\lfloor m/2\rfloor)$ cells. The move is
\[
(b_2,b_1)\mapsto
\begin{cases}
(1,1)&\ell\mapsto\ell+1,\\
(1,0)&r\mapsto r-1,\\
(0,1)&r\mapsto r+1,\\
(0,0)&\ell\mapsto\ell-1.
\end{cases}
\]
Two of the four moves enlarge the window, which is what the theorem in Section~\ref{sec:thm}
excludes. Both comparisons are strict: an exactly zero autocorrelation counts as negative, and a
window mean exactly equal to the mean of its first half counts as $b_2=0$.

\paragraph{What the $O(n)$ claims are, and are not.} Every read count in this subsection is a
construction with a measured step count: the simulator costs $O(1)$ reads per step, and the number
of steps is measured on a finite grid rather than bounded. Only the monotone case has a bound on
the steps, and that is why it is the only one stated as a theorem. A rule whose run were longer
than the tapes we ran would cost more than the table says.

Of the rules below, two are cheap with a ratio that stays flat across the sizes we can run --- the
canonical rule, and the monotone family of Theorem~\ref{thm:monotone}. Four others have a simulator
whose ratio grows with $n$: the interior-cut rule, the far-end rule, the median rule and the rule
that ignores the window. We report their numbers and leave them unclassified, because a ratio that
grows is not evidence of a linear bound and we cannot tell it from a slow drift on the samples we
have.

\paragraph{Monotone rules: $2n$ reads.} If every move of the rule shrinks the window, then each cell
leaves the window at most once and a rule that updates its feature when a cell leaves is exact in $2n$
reads: one read when the cell enters, one when it leaves. This is Theorem~\ref{thm:monotone} and it
covers every rule whose move table has no enlarging move.

\paragraph{The canonical rule: $5.3\,n$ reads.} The rule just defined enlarges the window
on two of its four moves, so the theorem does not apply. It is cheap anyway, and the reason is where
its cuts fall: it never chooses a cut action, so a cut happens only when a move would leave the
region, and the cell removed is then the boundary cell on that side. Keeping accumulators for the
region as well as for the window --- the window's are copied from the region's when a cut restarts the
headers --- costs $O(1)$ per step, and the measured cost is $5.3\,n$ reads per tape, flat from $n=32$
to $n=1024$ (Table~\ref{tab:sim}). The construction is in Section~\ref{sec:cost}.

\paragraph{All twenty-four move tables: $4.0$ to $5.1\,n$ reads.} The bound does not depend on which move
each pair of bits selects, only on how far a move can go. Replacing the table by any of the twenty-four
permutations of the four moves leaves the argument intact; measured over $120$ tapes at $n=256$, reads
per tape lie between $3.96\,n$ and $5.05\,n$, with no window ever rebuilt.

\paragraph{Steps of one, two and three cells: $5.0$ to $5.4\,n$ reads.} The same holds when a move moves
the header by two or three cells: a cut still removes the boundary cell only, and the accumulators still
lose one term. Measured reads per tape are $5.0$ to $5.4\,n$ with no rebuild.

\paragraph{A rule that cuts inside the window: $10.6$ to $17.3\,n$ reads.} A rule may also cut where it likes
rather than at the boundary. Take the canonical bits and moves and add one cut action: when the window
contains three consecutive ones, cut at the first of them and keep the left part. The rule is still
memoryless --- the decision is a function of the window --- but the new window is not a deletion of the
old one, so the accumulators have to be read again from $L$ to the cut. Measured over $200$ tapes, reads
per tape are $17.0\,n$, $17.3\,n$ and $10.6\,n$ at $n=128,256,512$: the ratio does not grow with $n$, so
the rule is pointer, and the rebuilding it does is paid for by the cuts it avoids.

\paragraph{A rule that ignores the window: $0.7\,n$ steps.} With $b_2=[2m>n]$ and $b_1$ the sign of the
lag-one autocorrelation, the walk marches inward and meets on most tapes in $0.69$ steps per cell, which
is pointer by the definition. We mention it because a fitted model scores well on it while the rule is exactly simulable:
a pointer strategy can look statistically rich and still be computable exactly, which is why $R^2$ on a
pointer strategy says nothing about the ceiling.

\paragraph{A rule that cuts at the far end: between $5.6\,n$ and $17\,n$ reads, and not classified.} This rule needs an action that the machine
of Section~\ref{sec:machine} does not have --- a cut at an interior position $p$, discarding the
cells to its right --- and we include it here as an extension, because it is the natural test of
whether interior cuts are expensive. The canonical bits and moves, plus that action, applied when the
window is longer than $n/2$ and its second half contains a cell equal to one: cut at that cell and
keep the left part. The rule is memoryless and both of its bits are aggregates. An
earlier draft reported $\Theta(n^2)$ reads for it; that number was measured by an implementation that
rebuilt the window after every cut. A cut action at a boundary removes one cell from the side it
names, so an accumulator kept over the region survives it; an interior cut removes a suffix and is a
different action, treated below, and the window that a cut resets to the region can be copied from the
region's accumulators in constant time. Carrying the two sets, reads per tape on tapes that produce an
output are $5.55\,n$, $6.08\,n$, $6.92\,n$, $13.5\,n$ and $17.0\,n$ at $n=128,256,512,1024,2048$,
against $31.8\,n$, $68.4\,n$, $179.0\,n$, $1352\,n$ and $3186\,n$ for the rebuild-per-cut
implementation. The two agree on $160/160$ output-producing tapes at each size, so the difference is the
implementation and not the rule.

The ratio grows by a factor of three over a sixteenfold range, which the meeting counts alone
might explain --- fourteen tapes at $n=1024$ and ten at $n=2048$ --- but we cannot separate the two
effects. We therefore do not put this rule in either class: it is not evidence for the cheap side,
and it is not used in Section~\ref{sec:ceiling}.

\paragraph{The median rule: three counters, $8.9$ to $13.8\,n$ reads, and not classified.} Take the canonical table and cut
policy and replace $b_2$ by the comparison of the window mean with the window \emph{median}. A median
is not an aggregate over an alphabet that grows with the window, but this tape has three values: with
$c_0,c_1,c_2$ the counts of $0,\tfrac12,1$ in the window, the $k$-th smallest cell is $0$ if
$k\le c_0$, $1$ if $k\le c_0+c_1$ and $2$ otherwise, so the median --- and, for even $m$, the mean of
the two central order statistics --- is a function of three counters, each $O(\log n)$ bits.
Reads per tape on tapes that produce an output are $12.3\,n$, $8.9\,n$, $13.6\,n$ and $13.8\,n$ at
$n=128,256,512,1024$, against $550\,n$, $1{,}061\,n$ and $2{,}941\,n$ for a driver that recomputes the
median at every step; the two agree on $200/200$ tapes at $n=128,256,512$. The ratio grows by half over that range, as it does for the far-end rule, so we
do not put this rule in either class either: the construction is linear per step and
the trajectory is not.

\subsection{The walk class}
\label{sec:walk}

Four rules are in it. All keep the canonical move table and take their second bit from the
\emph{arrangement} of the window rather than from a sum; all four are simulated with $O(\log n)$ bits
of state; and in all four the cost is quadratic because the run is long, not because a step is
expensive.

\paragraph{The longest run: $0.62\,n^2$ reads from a long run, not from a costly bit.} Let $b_1$ be
the sign of ``the longest maximal run of cells equal to $1$ in the window exceeds $m/4$''. The
threshold is a constant fraction of the window, so the bit can be decided from a few witnesses: if
the window contains $s$ cells that are not $1$, the longest run exceeds $t=\lfloor m/4\rfloor$ exactly
when some gap between consecutive such cells, or between one of them and a boundary, exceeds $t$. A
greedy scan jumps $t+1$ cells at a time and looks back for a non-$1$ cell, so deciding $b_1$ reads
$5.69$ to $5.89$ cells per step and stores four positions, which is $O(\log n)$ bits. The total per
step is about eleven reads once the initial scan of the tape and one rebuild per cut are counted;
that is what makes $0.62\,n^2$ compatible with a mean of $.054\,n^2$ steps. Measured on the tapes that
produce an output, the cost is $0.62\,n^2$ reads at $n=256$ and $0.63\,n^2$ at $n=512$, against
$9.8\,n^2$ and $20.1\,n^2$ for a driver that recomputes the longest run at every step; the two agree
on $60/60$ tapes at each size.

What makes the rule expensive is the length of its trajectory. The number of steps scales as
a fitted exponent over $n=128,\dots,2048$ that depends on the estimator: $2.089\pm.026$ on the median
and $1.977\pm.042$ on the mean. We quote both because they disagree, and the mean ratio is the one the
cost statement uses. The mean of steps$/n^2$ is flat --- $.053$, $.052$ and
$.048$ at $n=1024$, $2048$ and $4096$ --- while the median falls over the same range, from $.034$ at $n=1024$ to
between $.018$ and $.024$ at $2048$ --- two runs of different size disagree there --- and $.0095$ at
$4096$, with the ninetieth percentile rising from $.135$ to $.165$. The quadratic cost is
therefore carried by the tail: a typical tape is cheaper than $n^2$, and the mean is what the
asymptotic statement is about.

\paragraph{Three more, with bits that cost one read per step.} The same table with $b_1$ replaced by
(a) the comparison of the number of maximal runs of equal cells with $m/2$, (b) the presence of a
fixed three-cell pattern, and (c) the test of whether the window is a palindrome. Each of these admits
a constant-cost update --- a count of adjacent unequal pairs, a count of matching start positions, a
pair of rolling hashes --- so a step costs one read. All three are nevertheless quadratic, because the
trajectory is: steps$/n^2$ is $.046$, $.051$ and $.054$ at $n=256$, and $.063$ to $.069$ at $n=512$
across the family.

\begin{center}
\small
\begin{tabular}{lrrrrr}
\toprule
rule & tapes & meetings & best $R^2$ at $n^2/\ln n$ & fold SE & control \\
\midrule
longest run & 8000 & 1919 & $.58$ & $.014$ & $-.02$ \\
fixed pattern & 500 & 152 & $.41$ & $.039$ & $-.11$ \\
palindrome & 2000 & 255 & $.37$ & $.050$ & $-.06$ \\
run count & 1800 & 330 & $.28$ & $.017$ & $-.07$ \\
\bottomrule
\end{tabular}
\end{center}

Four rules, one budget, and a spread of a factor of two in what a predictor can recover: the class
is not homogeneous in predictability, which is what a single number for $\kappa$ would hide.

At $n=512$ the pattern rule gives $.37$ on $65$ meetings and the other two fall to about $-.17$, on
$80$ and $43$ meetings. Those samples are too small to separate a real drop from noise, and the
shuffled controls are of the same size as the values, so we report the numbers and read no trend
into them.

Only the first of the four is evidence about a ceiling. The other three cost about $.05\,n^2$ reads
to simulate exactly, and the budget $n^2/\ln n$ is $.18\,n^2$ at $n=256$ and $.16\,n^2$ at $n=512$
--- larger than their exact cost. A predictor with that budget can therefore simulate them
outright and reach $R^2=1$, so their scores measure the gradient booster we fitted and not a
ceiling. The longest-run rule is different: its exact cost is $0.62\,n^2$, above the budget at every
size we can run, and it is the only one of the four we use in Section~\ref{sec:ceiling}. None of the four is predicted well, and the two with the lowest scores are not the two with
the widest laws --- the last column is the same to within a third across the four.

The four are not four unrelated objects. Counting the pair $(b_2,b_1)$ along the trajectory splits
them into two regimes: at $n=128$ the longest-run and palindrome rules have the pair $(1,0)$ on $92$ per cent of
steps and $(0,0)$ on about $8$, while for the run-count and pattern rules it is $(1,1)$ on about $90$.
So there are two dynamical regimes with two members each, and the class is a pair of regimes rather
than a list. Figure~\ref{fig:ws} shows the four at $n=128$, where the budget is small enough that two of them are
predicted not at all; at $n=256$, with the larger budget of the table above, all four land between
$.28$ and $.58$. The moral is the one an earlier draft got wrong. We had been classifying bits by
whether a sum can
carry them, and every candidate that failed that test turned out to have an $O(1)$-per-step
implementation. The property that separates the classes is not the cost of a step but the length of
the run: a rule whose bits rarely push the headers together keeps cutting and restarting, and
$\Theta(n)$ cuts with $\Theta(n)$ steps between them is $\Theta(n^2)$ reads whatever a step costs.

\paragraph{No shortcut is known.} Seven summaries of the tape computable in $O(n)$ reads were tested
against the exact stopping position: the argmax and argmin of the centred walk, the centre of mass of
the cells equal to $1$, the midpoint of the longest run of them, the median of their positions, a
block-sum argmax at scale $\sqrt n$, and the last index at which the running rate reaches $1/4$. None
reproduces the output: the best of them, the centre of mass, is exact on $6$ of $50$ tapes at $n=64$
and on $0$ of $58$ at $n=256$, with a mean absolute error of $.106\,n$ to $.127\,n$. Coarsening the
simulation to blocks of $2$, $4$ and $8$ cells does not degrade gracefully but collapses, returning
the right answer on $0$ to $3$ of $46$ tapes. And replacing $b_1$ by the constant false leaves the step
count at $.061$ to $.076\,n^2$, so the long trajectory is a property of the move table's dynamics and
not of the bit: with one bit forced constant, the headers can only shrink, every few steps cut, and
the run restarts.

\begin{figure}[p]
\centering
\includegraphics[width=.94\textwidth]{figs/ws_longrun.png}\\[1pt]
\includegraphics[width=.94\textwidth]{figs/ws_runs.png}\\[1pt]
\includegraphics[width=.94\textwidth]{figs/ws_pattern.png}\\[1pt]
\includegraphics[width=.94\textwidth]{figs/ws_pal.png}
\caption{The four walk rules at $n=128$, each simulated for as many steps as the budget $n^2/\ln n$
allows (about a third of the trajectory). Left panel of each pair: out-of-fold prediction against
the strategy's output, with the diagonal for reference. Right panel: the density of the strategy's
$\delta$ and of the predictor's, on the same bins. From top to bottom, the longest-run rule
($R^2=.53$), the run-count rule ($.02$), the fixed-pattern rule ($.46$) and the palindrome rule
($-.02$). The two rules that resist prediction have outputs of the same spread as the two that do not,
so the difference is not the width of the law.}
\label{fig:ws}
\end{figure}

\section{A theorem: monotone rules are predictable}\label{sec:thm}

\begin{theorem}[Monotone rules are predictable]\label{thm:monotone}
Fix a strategy $(\varphi,f)$ with $f$ monotone: every move is $\ell\mapsto\ell+1$ or $r\mapsto r-1$,
no cut action and no halt action is ever chosen, so
exactly one cell leaves the window per step. Suppose further that the predictor may read the initial window, and that the state $s_0$ determining
$\varphi$ on it is computable from that window in $O(n)$ reads with $O(\log n)$ bits of workspace ---
a hypothesis that a finite value set of $\varphi$ does not by itself supply. Suppose $\varphi$ is
\emph{incrementally computable}:
there is a machine $M$ with $O(\log n)$ bits of workspace which, given the window $[\ell,r]$, a state
$s$ that determines $\varphi(\tau[\ell..r])$, and the cell that is about to leave, produces in $O(1)$
steps a state $s'$ that determines $\varphi$ on the new window. Then there is a predictor with
$O(\log n)$ bits of workspace that computes $\delta$ in at most $2n$ reads.
\end{theorem}

\begin{proof}
The predictor keeps $(\ell,r)$ and $s$; initially it reads the whole window, $n$ reads, and computes
$s_0$ by the hypothesis just stated. At each step it evaluates $\varphi$ from $s$, applies $f$, and calls $M$ with the cell that
leaves. By induction on the step the predictor's $(s,\ell,r)$ equals the rule's at every step: the
base case is the initial window, and the step case is the definition of $M$. Hence the first time the
headers would coincide is the same for both, and the predicted $\delta$ equals the rule's. One cell
is read per step, so the total is $n$ for the initial window plus at most $n-1$ for the steps.
\end{proof}

The hypothesis is met by the two features of the canonical rule separately. The three window sums
$S_1=\sum a_i$, $S_2=\sum a_i^2$, $S_3=\sum a_ia_{i+1}$ are updated by subtracting one term, and
$\varphi$ can be evaluated from them; the mean of the first half is updated by moving its two ends by
at most one cell each. So both features are maintainable along a move, and a monotone rule that
reads them is cheap by the theorem.
This does not contradict the measurements of Section~5, and the difference is where the cost comes
from. A move changes the window by one cell, and there the maintenance is constant. A cut does not:
it restarts the window at the region, so the window's sums cannot be carried across it by a single
deletion. What keeps the canonical rule cheap is the second set of sums, the one kept over the
region: those lose one term at the cut, and the window's are copied from them
(Section~\ref{sec:cost}). The cost is then $O(1)$ per cut and not a rebuild, which is why the
canonical rule does not pay for the $n/2$ cuts it makes.

\paragraph{Scope.} The theorem covers monotone rules only. The canonical rule is not one: two of
its four moves enlarge the window. So the bound of $2n$ reads says nothing about it, and the
measurements of the next section are not a counterexample to the theorem but a statement about a
rule outside its hypothesis.

\section{What simulation costs}\label{sec:cost}

Cutting is an action the rule may take at any time, at any place; nothing in the model ties a cut to
an exit. The canonical rule never chooses it, and that is what makes its simulation cheap --- not the
machine. We record the observation for the canonical rule first, then say what remains true when a
rule cuts wherever it likes.

\paragraph{The canonical rule.} A cut fires only when a move would leave the region, so it removes the
boundary cell on the side the header would have left and the headers restart at $(\ell,r)=(L,R)$. Two
sets of accumulators therefore suffice: the region's, which lose one term at the cut, and the
window's, which are copied from the region's when the headers restart. Neither is ever rebuilt from
the tape, and the cost is $O(1)$ per step.

Carrying, in each set, the window sum, the sum of squares, the sum of adjacent products and a running sum for the head window $[\ell,\ell+k-1]$ and the two
boundary values $a_\ell$, $a_r$ gives a simulation with $O(1)$ reads per step. With cells held scaled
by two, both bits become exact integer comparisons, $Sk>S_hm$ for the mean test and $(m-1)R_r-AB>0$
for the sign of the autocorrelation. The state is a fixed number of integers, so $O(\log n)$ bits.

\begin{table}[h]
\centering
\small
\begin{tabular}{rccccc}
\toprule
$n$ & tapes & meetings & reads per tape & reads$/n$ & reads$/n^2$ \\
\midrule
32   & 300 & 85 & 168.2   & 5.26 & .164 \\
64   & 300 & 69 & 340.2   & 5.32 & .083 \\
128  & 300 & 74 & 685.2   & 5.35 & .042 \\
256  & 300 & 49 & 1361.8  & 5.32 & .021 \\
512  & 150 & 22 & 2733.0  & 5.34 & .010 \\
1024 & 100 & 20 & 5501.4  & 5.37 & .005 \\
\bottomrule
\end{tabular}
\caption{Reads per tape for the exact simulation of the canonical rule, on tapes that produce an
output. The ratio reads$/n$ is flat at $5.3$ across a thirtyfold range of $n$, so the fraction of the
budget $n^2$ that the simulation uses falls like $5.3/n$. Seeds and commands
are in the repository. All rows use the exact tie convention,
where an exactly zero autocorrelation counts as negative; emulating the published float convention instead
costs $6.0\,n$ reads, a $13$ per cent increase that changes no output.}
\label{tab:sim}
\end{table}

Two checks fix what ``exact'' means here. The simulation agrees with a from-scratch integer
recomputation of the same window, state and both bits, on every tape at $n=32,\dots,256$; the rows at
$n=512$ and $1024$ are not recomputed that way, since the from-scratch check is quadratic per tape.
It also reproduces the published floating-point implementation bit for bit when the same tie-break
is emulated, $250/250$ at $n=64$, $200/200$ at $n=128$ and $120/120$ at $n=256$; those samples are
smaller than the table's $300$ because the float emulation is the slower of the two drivers. The tie-break matters,
because $\mathrm{corr}>0$ is decided by roundoff at exactly zero autocorrelation, which happens on
$6$ to $14$ per cent of steps since the tapes are integer-valued. Under the exact convention, where a
tie counts as negative, the cost is $5.3n$ reads; under the published convention it is $6.0n$.

\paragraph{Cuts elsewhere.} A rule that cuts at an interior position, or cuts a block rather than a
single cell, restarts from a window that is not a deletion of the old one: it can be longer, and the
accumulators must be rebuilt from its cells. The cost is then $T+\sum_j w_j$. That was the source of
the quadratic measurement in an earlier draft, and it is not where the expense actually lies: a rule
can keep the region's accumulators and pay $O(1)$ per cut, which is how the far-end rule is simulated
in Section~\ref{sec:classes}. What makes the four rules of Section~\ref{sec:walk} expensive is
neither the cut nor the bit but the number of steps.

\paragraph{Where this sits in the classical picture.} Simulating a machine with several heads by one
with a single head is an old problem. A multi-tape or multi-head machine can be simulated by a
single-head one in quadratic time \citep{hartmanis-stearns1965}, by a two-tape machine in
$O(n\log n)$ \citep{hennie-stearns1966}, and Hennie's crossing-sequence argument gives an
$\Omega(n^2)$ lower bound for a single tape that recognises the palindromes
\citep{hennie1965}. The fastest possible simulation is therefore known to lie between a nonlinear
lower bound and the quadratic upper bound, and the gap has stayed open
\citep{kutrib2010}. Our measurements sit inside that gap: with $O(\log n)$ bits the canonical rule costs $5.3\,n$ reads,
and the four rules of Section~\ref{sec:walk} cost $\Theta(n^2)$. The other two points that are often
quoted for this trade-off --- $n^2/\log n$ at $O(\log^2 n)$ bits and $n$ at $O(n)$ bits --- come from
the checkpointing construction of an earlier draft; the first of them is an upper bound that
construction, not a measurement we make here, and the second is trivial, since $O(n)$ bits holds the
tape.

\paragraph{Multi-head automata and pebbles.} Two-way finite automata are no stronger than one-way
ones \citep{rabin-scott1959}, but adding heads does add power: $k+1$ heads are strictly stronger than
$k$ \citep{rosenberg1966}, and the hierarchy survives even when the heads carry no state
\citep{kutrib-malcher2009}. A machine that can tell when two of its heads coincide is a sensing-head
automaton \citep{petersen1995}, and a machine that may drop markers on the tape is a pebble automaton
\citep{blum-hewitt1967}. Our machine sits between the two: it has two heads, it can detect coincidence,
and it can cut --- but a cut is not a pebble. A pebble is placed where the machine chooses and can be
lifted; our two boundaries only move inward, the place where they stop is set by the first exit of the
walk rather than chosen, and a cut destroys the cells behind it for good. That is why the questions
asked here --- how many reads a given workspace buys, and how much of the output a given amount of
memory determines --- do not reduce to the known hierarchy results.

\section{What cannot be proved this way}\label{sec:barrier}

It is natural to conjecture the converse of the theorem: that a rule which may grow the window again
is not exactly predictable without running the walk. That conjecture is false, and the refutation is
structural rather than accidental.

\paragraph{A probe that cycles.} A rule of this shape is easy to write down and we give it without
counts: fix a word $w$ of length $L\approx\log_4 n$, advance $\ell$ while the window does not begin
with $w$, then bring $r$ to $\ell$, and send the run into a cycle when the position falls in a fixed
interval. Such a rule cycles on some tapes, has a spread output on the ones where it meets, and is
reproduced by a one-pass predictor on every tape where it outputs, because on those tapes the output
is a function of the first hit of $w$. So a rule can cycle, have a spread output and be perfectly
predictable \emph{conditional on producing one}; what the one-pass predictor does not do is say in
advance whether the run will output or cycle. We record the construction as an obstacle to any
classification by cycles, and not as a measurement.

\paragraph{What else fails.} Influence does not separate the classes: the exchange measurements of
Section~\ref{sec:lower} show that the canonical rule has a quarter to a half of its cells pivotal,
and a monotone rule that costs $2n$ reads has cells that matter just as much. Nor does cutting, for
the reason given above: a rule may cut and still be cheap.



\paragraph{Why a class-level lower bound is hard to state.} A lower bound of the form "every
rule in the class costs at least $g(n)$ reads" fails as soon as the class contains one cheap
rule, because that rule is a counterexample. Feature-syntax classes tend to contain a monotone
member, since a feature that can be updated in constant time is exactly what makes a monotone rule
monotone. So such bounds are hard to state for classes defined by the form of the feature, and the
two other natural conditions --- cycles, cutting --- are refuted by the probes above.
\section{The ceiling, and the rules that test it}\label{sec:ceiling}

Everything in Section~\ref{sec:cost} is about rules whose output can be computed in linear time. For
those, no statistical question arises: an exact simulation reaches $R^2=1$ and a fitted model's value
measures the model, not the rule. The conjecture is about the rules that are left.

\paragraph{Statistical conventions.} The ensemble is the one of Section~\ref{sec:machine}: two
independent uniform $\lceil n/2\rceil$-subsets of $[n]$, one per column, and the tape their average.
Write $\Delta_n\in[0,1]\cup\{\bot\}$ for the output random variable, with $\bot$ on the tapes that halt or repeat, and $P(\Delta_n\ne\bot)$ for the meeting rate. A run either outputs $\delta\in[0,1]$ or outputs nothing, and every $R^2$ in this paper is computed
\emph{conditional on producing an output}, on the tapes where $\delta$ exists; the number of such
tapes is reported beside each figure, and the meeting rate is the fraction of all tapes. We claim no
mode of convergence: every number here is an out-of-fold estimate at a stated $n$ and a stated
budget, and the conjecture is about the limit that those estimates would approach if the pattern held.

\begin{conjecture}[Ceiling]
Let a rule of this class have an output that cannot be computed in $O(n)$ cell reads with $O(\log n)$
bits of state. Then every predictor that reads the tape in $o(n^2)$ reads with $O(\log n)$ bits
of memory has $R^2\le\kappa$ for a constant $\kappa<1$.
\end{conjecture}
\paragraph{What the budget excludes.} The simulations of Section~\ref{sec:walk} cost $0.62\,n^2$ reads,
so a budget that is $o(n^2)$ is eventually below the cost of following a trajectory, while a budget of
$c\,n^2$ is not: a predictor that spends $c\,n^2$ can follow the run to its end and reach $R^2=1$. That
is why the conjecture names $o(n^2)$, and every measurement below stays inside it: the budgets used run from
$19\,n$ to $82\,n$, all of which are $o(n^2)$.

\paragraph{The class has four candidates.} Four rules are candidates for the hypothesis of the
conjecture --- no $O(n)$ simulator is known for any of them --- and they are listed
in Section~\ref{sec:walk}: the longest-run rule and the three whose bits cost one read per step but
whose trajectories still run for $\Theta(n^2)$ steps. On the first of them we can measure the quantity
the conjecture bounds.

\paragraph{The ceiling moves with the budget, so we report the curve.} On the longest-run rule at
$n=256$ it scores $.58$ at the $n^2/\ln n$ budget of Section~\ref{sec:walk}; raising the budget towards
$n^2$ but staying inside $o(n^2)$ buys more. With $1919$ meetings and a shuffled-label control between
$-.003$ and $-.026$, a predictor that knows the first $s$ steps of the trajectory scores

\begin{center}
\small
\begin{tabular}{rrrrrrrr}
\toprule
$s$ & 24 & 32 & 48 & 64 & 96 & 128 \\
reads & $18.2\,n$ & $22.9\,n$ & $30.8\,n$ & $37.0\,n$ & $45.4\,n$ & $49.6\,n$ \\
$R^2$ & $.416$ & $.442$ & $.486$ & $.529$ & $.579$ & $.580$ \\
\bottomrule
\end{tabular}
\end{center}

All rows come from $8000$ tapes at $n=256$, of which $1919$ produce an output. A second seed of
$8000$ tapes, with $1939$ meetings, gives $.583$ at the same budget, and the two agree to $.004$.
The bootstrap resamples the output-producing tapes with replacement, four thousand times; in each
resample it recomputes the out-of-fold $R^2$ at every budget of the grid from the stored predictions,
takes the maximum over the grid, and reports the fifth percentile of that maximum. Taking the maximum over the budget grid inside each bootstrap resample --- the selection the
curve actually performs --- the one-sided $95$ per cent bound is $.563$ for the first seed and
$.561$ for the second, so we record
\[
\kappa\ \ge\ .56\quad(\text{one-sided }95\%,\ \text{selection-adjusted}).
\]
The same selection-adjusted procedure at $n=512$, on $900$ tapes of which $198$ produce an output,
peaks at $.48$ with a bound of $.39$. The drop from $.56$ is about three standard errors of that
sample, so it is not a null result; but with $198$ meetings the interval is wide and a power
statement would need several thousand tapes. We read it as a check rather than as an estimate: the
value does not rise with $n$, whereas a ceiling of one would rise.
The selection over the budget grid is accounted for by maximising inside each resample, but nothing
in the construction accounts for a search over rules, over feature sets, or over the number of
trajectory prefixes, and a wider search would lower the bound again. Nor does it locate $\kappa$: a cheaper way
of evaluating the window bit buys more of the trajectory inside the same budget and pushes the estimate
up, as the difference between $5.7$ and $9$ reads per step already shows. What the curve does establish
is that the value is not an artefact of a starved predictor --- the largest budget buys about a third of
the exact simulation's reads and the best use we found of them leaves more than a third of the variance
unexplained --- and that where the curve saturates is a better question than the constant.

\section{The lower bound}
\label{sec:lower}

The classification of Section~\ref{sec:classes} puts a rule on the expensive side when no $O(n)$
simulator is known. Turning that into a theorem needs a lower bound. This section gives the structural
facts we can prove, each checked exhaustively on small $n$ --- and then shows that they
cannot settle the classification, for a reason that changes what the missing argument has to be.

\subsection*{The rule, restated}

Let $A_j=\mathrm{mean}(a[L..L+j-1])$ be the average of the first $j$ cells of the region, and let
$m=r-L+1$ be the window when the window starts at $L$.

\begin{lemma}[The second bit is a prefix-average test]\label{lem:prefix}
$b_2([L,r])=0$ if and only if $A_{\lfloor m/2\rfloor}\ge A_m$.
\end{lemma}

\begin{proof}
$b_2=1$ means the mean of the first $k=\max(1,\lfloor m/2\rfloor)$ cells is below the mean of the
window. For $m=2k$ and $m=2k+1$ both comparisons reduce, after multiplying by $2k$ and $m$, to
$A_k<A_m$.
\end{proof}

Checked on $20{,}000$ random windows of lengths $2$ to $200$: no mismatch. The sequence
$B_m=A_m-A_{\lfloor m/2\rfloor}$ is not monotone in $m$ --- it changes sign $9.7$, $14.9$ and $24.2$
times on average at $n=128,256,512$, and fewer than $4$ per cent of the pairs have at most one sign
change --- so the test cannot be located by binary search.

\begin{theorem}[The rule with $b_1$ forced false is a largest-concavity scan]\label{thm:concavity}
Set $b_1\equiv0$, so the move is $r\mapsto r-1$ when $b_2=1$ and $\ell\mapsto\ell-1$ otherwise.
Starting from the region $[L,R]$, let
\[
m^\star=\max\{m\in[2,R-L+1]\;:\;A_{\lfloor m/2\rfloor}\ge A_m\},
\]
with $m^\star$ undefined if the set is empty. Then: if $R-L+1<3$ the run halts; if $m^\star$ is
undefined the headers meet and the output is $L/(n-1)$; if $m^\star=R-L+1$ the region is cut on both
sides; otherwise it is cut on the left. One epoch therefore costs as many $b_2$ tests as the scan from
$R$ down to $m^\star$, about $0.1\,n$ on average.
\end{theorem}

\begin{proof}
An epoch starts with $\ell=L$ and $r=R$. If $b_2(L,R)=0$ the move is $\ell\mapsto\ell-1$, which
leaves the region, and since $r=R$ as well the cut is the double one. If $b_2(L,R)=1$ the move is
$r\mapsto r-1$, which stays inside and leaves $\ell=L$; the same test then applies to the shorter
window, and $r$ decreases by one at each step until either $b_2(L,r)=0$ --- and then the move
$\ell\mapsto\ell-1$ leaves the region with $r<R$, so the cut is the left one --- or $r=L$, at which
point $\ell=r$ and the run outputs $L$. The largest $m$ with $b_2(L,L+m-1)=0$ is therefore exactly
the value of $m^\star$ that decides the epoch, and the three cases of the statement correspond to
the three endings.
\end{proof}

The proposition was also checked, tape for tape against the naive implementation, on every admissible \emph{sequence} for
$n=6,\dots,13$ --- $141$ sequences at $n=6$ and $201{,}643$ at $n=13$, all agreeing. Since the walk
depends on the tape only through the sequence of cell values, this is exhaustive for the dynamics;
what it is not is exhaustive for the ensemble, which weights a sequence by the number of matrices
that realise it, and that weighting is what the measured rates in Section~\ref{sec:walk} carry. This is the structural reason the
trajectory is long: each epoch rescans the region from its right end down to the largest concavity of
the prefix averages, and there are $\Theta(n)$ epochs.

\begin{corollary}[The scan is quadratic]
For the rule with $b_1\equiv0$, a run performs at most $\binom{n+1}{2}$ tests of $b_2$, hence $O(n^2)$
steps.
\end{corollary}

\begin{proof}
Every epoch ends with a cut, and a cut removes at least one cell, so there are at most $n$ epochs. An
epoch on a region of length $M$ tests $b_2$ at most $M$ times, because the scan moves $r$ down by one
at each test. Summing over the epochs and using that the region only shrinks gives at most
$\sum_{M=1}^{n}M=n(n+1)/2$ tests.
\end{proof}

So the quadratic cost of the walk rules is not only measured: for the $b_1\equiv0$ rule it is proved as
an upper bound, and the measurements of Section~\ref{sec:walk} put the mean at $.06\,n^2$ with a fitted
exponent of about $2$. The matching lower bound is proved, unconditionally.
The passage from the $b_1\equiv0$ rule to the full one is not proved; Lemma~\ref{lem:short} says the difference is confined to windows of
length below $4\rho=O(\log n)$, and the measurements agree to $2.3$ per cent at $n=64$ and to
four decimal places from $n=256$ on.

\paragraph{A lower bound on the first epoch.} The first epoch is a statement about the whole tape,
and there the quadratic cost can be bounded below. Write the tape in doubled units,
$a_i=2\tau_i\in\{0,1,2\}$, so that the cell mean is $1$ and $s^2=\tfrac12$ --- the hypothesis
$s^2\ge\tfrac14$ used below is then satisfied with room to spare; in the original units of
Section~\ref{sec:machine} the same cell has variance $\tfrac18$. Let $X_m=\sum_{i\le m}(a_i-1)$, and put $\mathcal V_M(\theta)=(X_{\lfloor\theta M\rfloor}-2X_{\lfloor\theta M/2\rfloor})/(s\sqrt M)$.
Under the ensemble $\mathcal V_M$ converges in $C[0,1]$ to $\mathcal V(\theta)=B(\theta)-2B(\theta/2)$,
with $B$ a Brownian bridge, and the discrete test differs from $\mathcal V_M$ by $O(1/(s\sqrt M))$
uniformly, so for every $c\in(0,1)$
\[
\liminf_{M\to\infty}P(\mathrm{cost}_0\ge cM)\ \ge\ P\Big(\sup_{\theta\in[1-c,1]}\mathcal V(\theta)<0\Big)=:p(c).
\]
A Gaussian concentration bound for that supremum (Borell--Tsirelson--Ibragimov--Sudakov, with
Dudley's entropy estimate), applied conditionally on $\mathcal V(1)=b$ and then integrated over $b$ of
both signs --- which is where a factor $\tfrac12$ enters, and forgetting it is an error of a factor
two --- gives $p(.05)\ge.080$. The display may use $\sup\mathcal V<0$ because
$\mathcal V\stackrel{d}{=}-\mathcal V$, so that event has the same probability as
$\inf\mathcal V>0$. Hence
\[
E[\mathrm{cost}_0]\ \ge\ .004\,M
\]
for all large $M$. The measured constant is $.130$ at $n=4096$ over $3000$ tapes, so the theorem is
thirty times weaker than the truth; it is nevertheless a bound, and the only bound in the paper that
is linear rather than quadratic in the wrong direction.
The limit is attained to within one per cent in simulation: at $n=4096$ over $3000$ tapes, $P(\mathrm{cost}_0\ge cM)$ equals the
Brownian value to within one per cent at $c=.05,.1,.2,.3,.5$ (ratios $0.991$ to $1.003$), and the
sharpest constant the \emph{simulated} limit gives is $c\,p(c)=.057$ at $c=.3$; the rigorous constant
is the $.004$ above, and the gap between them is the price of the concentration bound, not of the
limit. A purely combinatorial proof of the
same statement, with the worse constant $3.7\times10^{-4}$ and no functional limit theorem, is in
the repository; it rests on two exact lemmas and a scale separation between the chord and the
bridge of a four-block decomposition.

\paragraph{The run costs $\Theta(n^2)$.} The passage from one epoch to the whole run cannot be made by
conditioning --- bounding each epoch given the history is exactly what fails pathwise --- so the epochs
are compared with suffixes instead. The frontier inequality says that for $R\le R^{\prime\prime}$ the
cost of a window $[L,R]$ is at least the cost of $[L,R^{\prime\prime}]$ minus $R^{\prime\prime}-R$.
The reason is not that a longer window costs more --- it need not --- but that the tests at every
$m\le M$ are the same for both windows, since they read only cells up to $L+m-1$; the scan of the
shorter window is therefore the restriction of the longer one's, and the two costs differ by at most
the number of added cells. Take $R^{\prime\prime}=n-1$ and write $\delta_i=n-1-R_i$. Since $L$
advances by one per epoch and $R$ by at most one, $\delta_0=0$ and
$\delta_{i+1}-\delta_i\in\{0,1\}$, so $\delta_i\le i$ and
\[
\mathrm{cost}_i\ \ge\ \mathrm{cost}(L_i,n-1)-i .
\]
Three things must be said about the next step. The object bounded is $\mathrm{cost}_i\cdot 1\{A_i\}$,
not a conditional expectation: the event $A_i$ that epoch $i$ exists is not a function of the prefix,
since the first epoch tests the whole tape, so the suffix is conditioned on the history and no
unconditional law applies to it directly. What is used instead is
\[
E[\mathrm{cost}_i 1\{A_i\}]\ \ge\ E[\mathrm{cost}(L_i,n-1)]-(n-i)P(A_i^c)-i,
\]
the middle term being the trivial bound $\mathrm{cost}\le n-i$ on the complement: the conditioning is
paid for, not assumed away.

Second, a suffix does not have balanced counts, and the single-epoch theorem was stated for a balanced
multiset. The extension is true and is proved. The test is exactly invariant under adding a constant to
every cell, so the multiset may be centred before it is normalised; the covariance identity
$\mathrm{Cov}(Y_i,Y_j)=s^2i(M-j)/(M-1)$ then holds for \emph{any} multiset, and the contrast of the
resulting bridge is the same functional $\mathcal V$; and a quantitative CLT for linear permutation
statistics, proved by Stein's method of exchangeable pairs, has an error
$2.89\,(K/s)/\sqrt m+O(1/M)$ that depends on $K/s\le4$ alone and not on the composition. The
hypothesis $s^2\ge1/4$ fails only on suffixes that are almost constant, which happens with
probability $e^{-\Theta(M)}$.

Third, $q_N=P(A_N^c)$, the probability that the run meets before epoch $N$. A meet forces, for some
$i<N$, the absence of a concavity at every scale of the fixed block $[i,i+n-2i-1]$, because the test
set is an initial segment in $m$; hence $P(\text{meet before }N)\le\sum_{i<N}f(n-2i)$, where $f(M)$ is
the probability that a \emph{fresh} uniform $M$-window has no concavity at all. On the dyadic chain
$X_{2^j}>2X_{2^{j-1}}$ the normalised walk is a supermartingale with
$E[e^{-3u_{j+1}}\mid\text{survival}]\le\frac14e^{-3u_j}$, which gives $f(M)\le2100\,M^{-2}$; summing
over the $\Theta(n)$ blocks, $P(\text{meet before }(n-1)/4)\le1050/n$.

Summing the epoch bound over the first $N=c_0n/(1+c_0)$ epochs, then,
\[
E[T]\ \ge\ \Big(\frac{c_0^2}{2(1+c_0)}-o(1)\Big)\,n^2 ,
\]
with $c_0=.004$ from the rigorous half of the previous paragraph, giving $8.1\times10^{-6}n^2$, and
with the measured $c_0=.130$ giving $7.5\times10^{-3}n^2$. Against the upper bound $\binom{n+1}{2}$
the run is quadratic in expectation, with no hypothesis left and a gap of five orders of magnitude
between the two constants.

\paragraph{Quadratic in expectation, not on every tape.} The corollary bounds the cost; the
measurements say what it is. Over $2000$ tapes at $n=1024$, $E[T]/n^2=.0540\pm.0014$, and the ratio
stays between $.048$ and $.059$ from $n=1024$ to $4096$. It is not a pathwise statement: the median
is $.025\,n^2$ against a mean of $.054\,n^2$, the standard deviation is of the order of the mean, and
$P(T\le n)=.17$ --- on those tapes almost every epoch is a double cut and the run costs $0.6\,n$. So
"the walk asks $\Theta(n^2)$ questions" is an expectation and not a law.

\paragraph{What controls the cost.} One statistic carries most of it. Write $\Delta$ for the
difference between the sums of the two halves of the tape. Then $E[T/n^2\mid\Delta\ge0]=.090$ to
$.104$ and $E[T/n^2\mid\Delta<0]=.012$ to $.016$, stable over $n=128$ to $4096$; the correlation of
$T/n^2$ with $\Delta/\sigma_\Delta$ is $+.74$, and $88$ per cent of $E[T]$ comes from the
$\Delta\ge0$ half. The drift persists into the sub-tapes --- $52$ per cent of epochs meet a window
with positive drift against $21$ per cent of fresh windows --- which is why the cost does not
concentrate.

\begin{lemma}[$b_1$ fires only on short windows]\label{lem:short}
Let $\rho$ be the longest run of cells equal to $1$ in the whole tape. If $b_1=1$ at a window of
length $m$ then $m<4\rho$. Since $\rho=O(\log n)$ with high probability, every window longer than
$O(\log n)$ is governed by Theorem~\ref{thm:concavity}.
\end{lemma}

\begin{proof}
If $b_1=1$ the window contains a run of more than $m/4$ cells equal to $1$; that run is a run of the
whole tape, so $\rho>m/4$ and $m<4\rho$.
\end{proof}

Measured over $40$ tapes per size, no tape violated $m<4\rho$, and $b_1$ was true on $0.84$, $0.31$,
$0.12$ and $0.03$ per cent of steps at $n=64,128,256,512$ --- about five steps per tape, always at the
tail of an epoch. Forcing $b_1\equiv0$ leaves steps$/n^2$ at $.0673$, $.0484$, $.0645$, $.0688$, against
$.0658$, $.0482$, $.0645$, $.0688$ for the full rule.

\begin{theorem}[General upper bound]\label{thm:upper}
Every strategy of the model is simulable with $O(n^2)$ bits of state and $O(n^3)$ steps, where a step
is one evaluation of $\varphi$ on the current window at unit cost.
\end{theorem}

\begin{proof}
The region only shrinks, so it changes at most $n$ times. Between two changes $[L,R]$ is fixed and the
headers range over the $\binom{n+1}{2}$ pairs with $L\le\ell\le r\le R$; the rule is deterministic, so
the segment either meets or revisits a pair, and a revisit ends the run. Hence a segment has $O(n^2)$
steps and a run $O(n^3)$, and detecting a revisit needs only the pairs of the current segment because
the region never grows.
\end{proof}

So the cost of a rule in this model lies between the certificate bound below and $O(n^3)$ steps, and
the four rules of Section~\ref{sec:walk} sit at $\Theta(n^2)$ measured, with a $0.62\,n^2$
construction.

\subsection*{Why none of this settles the classification}

\paragraph{Why memoising does not help.} A natural objection is that the tests repeat: the scan asks
the same question of shifting windows, so a simulator could cache the answers. It cannot. Every
contrast $(L,m)$ --- the comparison of the mean of $a[L..L+\lfloor m/2\rfloor-1]$ with that of
$a[L..L+m-1]$ --- is asked at most once in a run, because the anchor $L$ increases by exactly one per
epoch and the scan within an epoch uses a fixed anchor. What is reusable is the prefix sums and not
the decisions, and one left-to-right pass with three counters evaluates a whole epoch in
$O(\log n)$ bits. The run is a single-anchor sweep; all of its cost is the iteration, and none of it
is redundant computation. The one further model in which a superlinear bound does follow is a machine
whose memory is additive and whose decisions are exact: there the same accounting gives
$T\ge\sum_e\max(0,\ell_e-2r)$ over the epoch lengths $\ell_e$, which at $s=O(\log n)$ is again
$\Omega(n^2)$. Dropping additivity is the branching-program problem, and dropping exactness returns
the certificate bound $n-1$. A machine that reads in a fixed order in $p$ passes does no better: it
crosses each cut $2p$ times, so the cut lemma gives $p=\Omega(n/s)$ and $T=\Omega(n^2/s)$, and it
cannot escape by choosing a cheaper order, because the one-way communication complexity of the output
is the same for every balanced cut --- about $0.73\,n$ bits, measured exhaustively at $n=6,8,10,12$.

\paragraph{Why no read-counting argument can separate.} The exchange argument is the general form of
every read-based bound: a simulator that reads a set $S$ of cells behaves identically on all tapes
agreeing with the input on $S$, so $T\ge|S|$, and $|S|\le n$ by definition. The bound therefore caps
at $n$, and $n$ is on the cheap side of Definition~\ref{sec:classes}. It can be sharpened: computing minimum
certificates exhaustively --- the smallest set of cells that determines the output --- gives
$K_{\mathrm{global}}=n-1$ for $n=6,\dots,13$: every set of $n-2$ cells fails to determine the output
somewhere, and $n-1$ always suffices. An earlier run reported the opposite, and the cause was a
projection written over the cells outside the read set rather than inside it, which silently tests
two-cell sets; with the direction corrected the result reproduces independently. What survives is the structural point, which does not depend on it: any
argument that counts the cells a machine reads can prove at most a linear bound, and a linear bound
is exactly what the cheap class allows. \emph{The classification is therefore not a statement about
distinct cells at all; it is a statement about repeats.} It cannot be a statement about space
either, since $\delta$ is computable with $\Theta(\log n)$ bits. What has to be proved is that no
machine with $O(\log n)$ bits computes the output in $O(n)$ reads, which --- since a read set never
has more than $n$ cells --- means proving $\Omega(n\,\mathrm{polylog}\,n)$ repeat reads.

\paragraph{An $O(n)$-read algorithm exists and is wrong on half the tapes.} The scan of
Theorem~\ref{thm:concavity} can be shortened by assuming that the largest concavity position does not
increase when $L$ advances, which would make the whole run cost $O(n)$ tests. The assumption fails on about
a third of consecutive pairs at $n=12$, and the resulting algorithm is exact on
$92$, $86$, $54$ and $81$ per cent of the tapes at $n=8,10,11,12$ and on $52$ and $50$ per cent at
$n=64$ and $256$ --- with a mismatch rate that does not fall as $n$ grows. This is the closest thing to
a linear-read simulator that we have, and it is not one.

\paragraph{Why the obvious arguments cannot work.} Every lower bound of cut type --- a single cut,
all cuts aggregated, $p$-pass, pointer, external memory --- obeys one conservation law. Let $T$ reads
have mean jump length $\bar J$ on a machine with $s$ bits, and let $\mathrm{CC}_k(\delta)$ be the
two-way communication complexity of $\delta$ at the cut after cell $k$. Then
\[
T\,\bar J\,s\ \ge\ \sum_{k=1}^{n-1}\mathrm{CC}_k(\delta)-ns .
\]
Aggregating over the $n$ cuts gains a factor $n$; charging the jumps loses exactly a factor $n$,
because a single read may cross all $n$ cuts. With $\sum_k\mathrm{CC}_k\le\log_2 3\cdot n^2/4$ this
caps the whole family at $.396\,n/\log_2 n$ reads, whatever the memory, and the conjecture is
stronger than the law by about $\log^2 n$. The two regimes are worth separating. If the reads
are sequential, $\bar J=1$, the law reads $T\ge\big(\sum_k\mathrm{CC}_k(\delta)-ns\big)/s$. What is
proved about the sum is an \emph{upper} bound, $\log_2 3\cdot n^2/4=.396\,n^2$, and an upper bound is
the wrong direction here: to bound $T$ below one needs $\sum_k\mathrm{CC}_k$ to be large. It is:
measured over all balanced cuts the sum is $95$ to $98$ per cent of that maximum, and on that
measurement the single-head bound is $\approx.38\,n^2/s$, i.e. $\Omega(n^2/\log n)$ at
$s=O(\log n)$. That is a superlinear bound for single-head machines, but it rests on the measured size
of $\sum_k\mathrm{CC}_k$ and not on a proof of it, and we flag it as such. For random access the same
expression degrades to $\approx.38\,n/\log_2 n$, which is below the certificate bound $n-1$ and proves
nothing. What survives is a necessary condition on any counterexample: an
$O(\log n)$-bit machine that reads $n\,\mathrm{polylog}\,n$ cells must have mean jump
$\Omega(n/\log^3 n)$. The conjecture is a statement about machines that teleport. There is also no room for a space lower bound: $\delta$
is computable with $\Theta(\log n)$ bits, since the walk can be simulated, so $\omega(\log n)$ space
is not merely unproved but false. The conjecture is therefore a statement about repeat reads alone.

\paragraph{What is left open, stated precisely.} A machine with $s$ bits that halts cannot repeat a
configuration, so $T\le2^{s+O(1)}$: state caps reads, and $O(\log n)$ bits already allows
$\mathrm{poly}(n)$ of them. Twelve bits above $\log_2 n$ suffice for $n\,\mathrm{polylog}\,n$. The
question is therefore not a coarse separation of space from time but the shape of one curve: whether
$T(s)$ is a step at $s=\Theta(n)$ or the smooth $2^{\Theta(s)}$. Since $T=|S|+\text{repeats}$ with
$|S|\le n$, a proof needs $\Omega(n\,\mathrm{polylog}\,n)$ \emph{repeat} reads, which is a
branching-program lower bound of width $\mathrm{poly}(n)$ --- and no superlinear branching-program
bound is known for any explicit function. The three routes we tried each have a provable ceiling:
a one-way cut gives at most $0.79\,n/\log_2n$, counting states bounds $T$ from above rather than
below, and an adversary that keeps two outputs alive dies once $n-1$ cells have been read, because
every $(n-1)$-set is a certificate. The obstruction is the machine's memory and not the data's redundancy: $n$ reads always suffice, since
a machine with $O(n)$ bits can hold the tape, and the question is what $O(\log n)$ bits can do. We do not have
such an argument, and the two standard routes fail on this function: the prefix-separation test needed
for a one-way streaming bound is false (same-profile prefixes with identical output on every common
suffix occur in $45$ per cent of pairs at $n=6$ and $14$ per cent at $n=10$), and no fooling set for
the dynamical stopping position is known to us.

\section{The questions this leaves}

The question an earlier draft left open --- is the output of the \emph{canonical} rule computable in
$O(n)$ reads with $O(\log n)$ bits? --- is answered by a construction whose step count is measured rather than
bounded: $5.3\,n$ reads, from the observation that a cut removes a boundary cell and never a
window. That answer is about one rule and not about the class.
Four rules are in it; no $O(n)$ simulator is known for any of them, and the difference between them and the
canonical rule is the length of the trajectory rather than the cost of a step.

Three questions remain. The first is a lower bound for the walk class: we know no $O(n)$ simulator for
the four rules of Section~\ref{sec:walk}, and we have an argument that their trajectories are
$\Theta(n^2)$ steps long, but an argument is not a proof, and no reduction shows that following the
trajectory is necessary. The second is the constant: $\kappa\ge.56$ on one rule, and no
upper bound at all --- we do not know whether a predictor inside the budget can reach $.7$ or $.9$ on
the longest-run rule, and the number should be expected to move as better ways are found to evaluate
the window bit, since a cheaper evaluation buys more steps of trajectory inside the same budget. The
third is the one the classification does not reach: for an observer restricted to aggregates of the
initial tape, with no access to the trajectory, what is the best possible $R^2$? A formal statement
would need the restriction written as a class of computable functions rather than as a read budget,
and we do not have such a class.

\section*{Reproducibility}

Every number in this paper comes from a script in the project directory; the repository is
\texttt{github.com/peuaza/triangular-window-walk} and the machine-paper scripts are under
\texttt{bthm0/}, \texttt{probe\_ologn/} and \texttt{probe\_ologn2/}. They run with Python 3 and numpy 2.5.2,
scipy 1.18.1, scikit-learn 1.9.1 and matplotlib 3.11.2. Tapes are drawn with
\texttt{numpy.random.default\_rng(2026)} unless a paragraph says otherwise, each column a uniform
$\lceil n/2\rceil$-subset. The predictor is a gradient booster with 80 iterations, learning rate
$.05$, depth $3$, leaf floor $20$ and $\ell_2$ penalty $1$, on five folds; every reported $R^2$ is
out of fold with a shuffled-label control from the identical pipeline, and the fold SE is the
standard error over those five folds.

Table~\ref{tab:sim} is produced by \texttt{code/bthm0/\_bz\_base.py} and reproduced against the
published implementation by \texttt{\_bz\_canon2.py}; the four walk rules by
\texttt{\_arr\_longrun\_fast.py} (costs and agreement) and \texttt{\_phase\_arr2.py} (the budget
curve and the stored out-of-fold predictions); the other three walk rules by
\texttt{\_phase\_rule.py}; the bootstrap by \texttt{\_boot\_sel2.py}; the exchange measurements by
\texttt{\_flip\_canon2.py}; and the figures by \texttt{\_fig\_walk.py}. Everything in
Section~\ref{sec:lower} is in \texttt{probe\_ologn2/}, with the tables and commands in its
\texttt{REPORT\_EPOCH\_BOUND.md}, \texttt{REPORT\_LEMMA\_Q.md} and \texttt{REPORT\_MEET.md}; the
communication and certificate work is in \texttt{probe\_ologn/}.
\texttt{REPORT.md}.

\bibliographystyle{plainnat}
\begin{thebibliography}{9}
\bibitem[Luo(2026a)]{luo-tww} H.~Luo, \emph{A triangular window walk on random binary matrices: a
conjectured arcsine limit}, manuscript (2026).
\bibitem[Luo(2026b)]{luo-family} H.~Luo, \emph{Memory and randomness control the shape: a family of
limit laws from a triangular-window walk}, manuscript (2026).
\bibitem[Rabin and Scott(1959)]{rabin-scott1959} M.~O.~Rabin, D.~Scott, \emph{Finite automata and their
decision problems}, IBM Journal of Research and Development 3 (1959) 114--125.
\bibitem[Rosenberg(1966)]{rosenberg1966} A.~L.~Rosenberg, \emph{On multi-head finite automata}, IBM
Journal of Research and Development 10 (1966) 388--394.
\bibitem[Kutrib and Malcher(2009)]{kutrib-malcher2009} M.~Kutrib, A.~Malcher, \emph{On stateless
multihead automata: hierarchies and the emptiness problem}, Theoretical Computer Science 410 (2009)
1711--1721.
\bibitem[Blum and Hewitt(1967)]{blum-hewitt1967} M.~Blum, C.~Hewitt, \emph{Automata on a two-dimensional
tape}, Proc. IEEE Symposium on Switching and Automata Theory (1967) 155--160.
\bibitem[Petersen(1995)]{petersen1995} H.~Petersen, \emph{Automata with sensing heads}, Proc. ISTCS
1995, 150--157.
\bibitem[Hartmanis and Stearns(1965)]{hartmanis-stearns1965} J.~Hartmanis, R.~E.~Stearns, \emph{On the
computational complexity of algorithms}, Transactions of the AMS 117 (1965) 285--306.
\bibitem[Hennie and Stearns(1966)]{hennie-stearns1966} F.~C.~Hennie, R.~E.~Stearns, \emph{Two-tape simulation
of multitape Turing machines}, Journal of the ACM 13 (1966) 533--546.
\bibitem[Hennie(1965)]{hennie1965} F.~C.~Hennie, \emph{One-tape, off-line Turing machine computations},
IEEE Transactions on Electronic Computers 14 (1965) 35--47.
\bibitem[Kutrib(2010)]{kutrib2010} M.~Kutrib, \emph{Complexity of multi-head finite automata: origins and
directions}, Theoretical Computer Science 412 (2011) 2--36.
\end{thebibliography}

\end{document}
